\chapter{Methodology}
\label{chap:methodology}

This chapter presents the complete analytical workflow: the data used and the role each dataset plays, the physics-based degradation pipeline and its calibration, the protocol for measuring the domain gap, the domain-adaptation mechanism adopted, the cross-architecture distillation procedure, and the evaluation protocol.
It closes with a dedicated account of the tooling and reproducibility practices used throughout.

The data underlying this work is split into two categories.
The first is the co-registered triplet dataset itself: real $\Phi$-sat-2 acquisitions, obtained via the ESA Insula Perception platform, matched to concurrent Sentinel-2 scenes and to a physics-simulated $\Phi$-sat-2 rendering of the same Sentinel-2 scene, co-registered using the LightGlue deep feature-matching algorithm \cite{Lindenberger_2023_ICCV} and filtered for spatio-temporal proximity and cloud cover.
This dataset carries land-cover labels derived from ESA WorldCover 2021, providing the only currently available real-sensor ground truth.
The second category is simulated: a physics-degraded, graded version of the Sen1Floods flood-segmentation benchmark, produced with the same degradation pipeline, standing in for a real labeled $\Phi$-sat-2 flood dataset that does not yet exist because the mission has not captured an annotated flood event.

The overall workflow followed a prototyping-implementation-refinement logic rather than a fixed plan executed once.
An initial conceptual model (fine-tune TerraMind on real $\Phi$-sat-2 labels, measure the gap, close it) was implemented, evaluated against intermediate diagnostics (embedding visualizations, similarity metrics), and revised twice on the basis of what those diagnostics showed: first when a normalization defect was found to be corrupting an early performance baseline (Section~\ref{sec:normalization}), and second when a first domain-adaptation mechanism was found, empirically, not to work (Section~\ref{sec:da-negative-result}), motivating the mechanism actually adopted (Section~\ref{sec:contrastive-mechanism}).

\section{Data}
\label{sec:data}

\subsection{The co-registered triplet dataset}
\label{sec:triplets}

The main dataset is a collection of spatially and temporally co-registered image triplets, produced in collaboration with the ESA $\Phi$-lab: real $\Phi$-sat-2 acquisitions (obtained via the ESA Insula Perception platform), the corresponding real Sentinel-2B scene, and a physics-simulated $\Phi$-sat-2 rendering of that same Sentinel-2 scene (Section~\ref{sec:degradation-pipeline}).
Candidate scenes were filtered by spatio-temporal proximity between the $\Phi$-sat-2 and Sentinel-2 overpasses and by a cloud-cover threshold on both acquisitions, and precise pixel-level alignment was achieved using the LightGlue deep feature-matching algorithm \cite{Lindenberger_2023_ICCV}, which identifies corresponding semantic keypoints across the two sensors despite their differing spatial and radiometric characteristics.

The choice of Sentinel-2 as the paired source sensor follows from the two instruments' band layouts.
Figure~\ref{fig:srf} compares their spectral response functions and their per-band spatial resolutions.
Across 400--900\,nm each $\Phi$-sat-2 multispectral band has a Sentinel-2 counterpart at a comparable centre wavelength, and it is that correspondence which makes Sentinel-2 usable both as the paired view in the triplets and as the source the degradation pipeline (Section~\ref{sec:degradation-pipeline}) projects into the $\Phi$-sat-2 domain.
The correspondence is close but not exact and $\Phi$-sat-2 additionally carries a panchromatic band with no Sentinel-2 equivalent.
Spatially, $\Phi$-sat-2 samples every band at a uniform 4.75\,m, whereas Sentinel-2 provides 10\,m for B2, B3, B4 and B8 and 20\,m for B5, B6 and B7.
The spatial resolution mismatch is a factor taken into account when interpreting the domain gap in Chapter~\ref{chap:results-discussion}: the target sensor is the finer of the two, so a resolution mismatch cannot by itself account for a loss of performance on $\Phi$-sat-2, whereas the spectral differences visible in panel~(a) remain a candidate explanation.

\begin{figure}[htp]
    \centering
    \includegraphics[width=12cm]{fig_srf_s2_phisat2.png}
    \caption{Band correspondence between Sentinel-2B and $\Phi$-sat-2, after Long\'ep\'e et al.~\cite{longepe2024orbitalai}.
        (a) Normalized spectral response functions, Sentinel-2B above and $\Phi$-sat-2 below.
        (b) Per-band spatial resolution.
        The correspondence in (a) makes Sentinel-2 usable as the paired source; the residual differences in both panels are the sensor gap this thesis aims to measure.
    }
    \label{fig:srf}
\end{figure}

This dataset carries labels for a single task: land-use/land-cover (LULC) classification, derived from ESA WorldCover 2021 and rasterized onto each triplet's real-$\Phi$-sat-2 footprint.
Its role in this thesis is twofold: it is the only real-sensor ground truth available for any downstream task, and its co-registration is what makes it possible to measure sensor-induced covariate shift in TerraMind's latent space (Section~\ref{sec:gap-measurement}) without the scene-content and geography confounds present in prior domain-gap studies that compare unrelated source and target datasets.

\begin{figure}[htp]
    \centering
    \includegraphics[width=12cm]{fig_patch_examples.png}
    \caption{Four representative triplets, each row showing the same scene as observed by real $\Phi$-sat-2, the physics-simulated $\Phi$-sat-2 rendering, and Sentinel-2B (each channel independently contrast-stretched for display).
        The real acquisitions show visibly grainier, noisier texture than either the simulated or source Sentinel-2 view, most apparent in the dark water/shadow region of patch 171127, consistent with the sensor noise characteristics the degradation pipeline (Section~\ref{sec:degradation-pipeline}) is designed to reproduce.
    }
    \label{fig:patch-examples}
\end{figure}

Figure~\ref{fig:patch-examples} shows four such triplets.
Qualitatively, the real acquisitions are consistently grainier than their simulated or source counterparts, and occasionally differ subtly in overall tone; the simulated column is, by construction, closer in appearance to Sentinel-2 than to the real sensor it is meant to approximate, a point returned to below.

\subsubsection{Scale and coverage}

After discarding a small number of corrupted source products, the dataset comprises 253{,}228 co-registered triplet patches drawn from 1{,}299 distinct $\Phi$-sat-2 products (individual acquisitions), a median of 209 patches per product (up to 713 for the largest).
Acquisitions span 2024-08-21 to 2026-04-15 (21 months), with collection volume markedly higher from mid-2025 onward than in the initial ramp-up period.
Real-$\Phi$-sat-2 and Sentinel-2 acquisition dates are matched to within a median of 6.8 days and a maximum of 14.9 days, not simultaneously, which is an acceptable temporal offset for a largely stable target such as land cover but a caveat worth stating rather than assuming away.
73.1\% of retained patches have under 5\% total cloud cover, and only 0.199\% of pixels are unlabelled after the ESA WorldCover rasterization step.

\begin{figure}[htp]
    \centering
    \includegraphics[width=\textwidth]{fig_acquisition.png}
    \caption{Left: distribution of the absolute time lag between matched $\Phi$-sat-2 and Sentinel-2 acquisitions (median 6.8 days, maximum 14.9 days, $n=253{,}228$).
    Right: monthly acquisition volume, Aug.
    \ 2024--Apr.\ 2026, showing the collection ramp-up in the first eight months.}
    \label{fig:acquisition}
\end{figure}

Geographically, the 1{,}299 source products span latitudes $-54^{\circ}$ to $76^{\circ}$ with global but uneven coverage, concentrated over western North America, the southern Andes, the Mediterranean basin and Middle East, and parts of East and South Asia (Figure~\ref{fig:geography}).
By K\"oppen--Geiger climate group, the patch population is dominated by arid (41\%) and temperate (29\%) climates, with tropical (14\%), continental (10\%), and polar (1\%) regions comparatively under-represented (4\% unclassified).

\begin{figure}[htp]
    \centering
    \includegraphics[width=\textwidth]{fig_geography.png}
    \caption{Left: spatial density of $\Phi$-sat-2 products (hexbin, log-scaled patch count per cell).
        Right: distribution of patches by K\"oppen--Geiger climate group\cite{beck2018present}.
    }
    \label{fig:geography}
\end{figure}

The eleven ESA WorldCover classes present are similarly imbalanced by construction rather than by sampling choice: grassland (20.8\%), tree cover (20.0\%), bare/sparse vegetation (16.2\%), built-up (14.4\%), and cropland (13.4\%) account for the large majority of labelled pixels, while snow/ice, mangroves, and moss/lichen each fall below 0.6\% (Figure~\ref{fig:landcover}).
The typical patch is fairly homogeneous: a single class occupies a median of 73\% of any given patch's labelled area.

\begin{figure}[htp]
    \centering
    \includegraphics[width=\textwidth]{fig_landcover.png}
    \caption{Left: share of all labelled pixels by ESA WorldCover class.
        Right: number of patches in which each class is the dominant (plurality) class.
    }
    \label{fig:landcover}
\end{figure}

\subsubsection{Train/validation/test split}

The dataset is split 80/10/10 (202{,}582 / 25{,}323 / 25{,}323 patches) at the patch level with a fixed seed, and the resulting class-prior distribution is stable across all three splits, confirming the split does not introduce label-distribution shift (Figure~\ref{fig:splits}, left).
However this introduces a different problem, because splitting is patch-level rather than product-level, all 1{,}299 source products are represented in the training set, and every one of the 1{,}287 distinct products contributing to the test set also contributes patches to training (Figure~\ref{fig:splits}, right).
Since patches from the same product are spatially adjacent or overlapping crops of one acquisition, this means the test set is not a clean geographic or scene-level hold-out, and real-data performance figures reported in Chapter~\ref{chap:results-discussion} should be read as an optimistic upper bound on generalization to genuinely unseen scenes rather than as a strict out-of-sample estimate.
A product-level (rather than patch-level) re-split is the direct remedy and is noted as future work in Section~\ref{sec:discussion}.

\begin{figure}[htp]
    \centering
    \includegraphics[width=\textwidth]{fig_splits.png}
    \caption{Left: class prior is stable across the train/validation/test split.
        Right: every one of the 1,287 products contributing to the test set also contributes patches to the training set: the split is patch-level, not scene-level.
    }
    \label{fig:splits}
\end{figure}

\subsubsection{Measured sensor discrepancy}

Three independent measurements quantify the sensor discrepancy this dataset is designed to expose, prior to any model being involved.

First, the raw radiometric scales of the two sensors heavily differ (Figure~\ref{fig:band-boxplots}): real $\Phi$-sat-2 digital numbers have per-band medians of about 55--150 (12-bit range, with measurable saturation of 2.1--2.2\% of pixels in the RE1/RE2 bands and dropout of 0.3--0.36\% in RE2/RE3), against Sentinel-2 and simulated-$\Phi$-sat-2 medians around 2{,}200--3{,}300, consistent with the Level-1C (L1C) Top-of-Atmosphere (TOA) reflectance-scaled DN convention.
The simulated and raw Sentinel-2 domains are close to statistically indistinguishable at this level (their boxes are overlapping in Figure~\ref{fig:band-boxplots}), confirming that the physics-based degradation pipeline (Section~\ref{sec:degradation-pipeline}) perturbs fine spatial and spectral structure rather than intensity statistics.
The domain gap this thesis is concerned with is therefore a real-vs-simulated gap, not a simulated-vs-source gap.

\begin{figure}[htp]
    \centering
    \includegraphics[width=\textwidth]{fig_band_boxplots.png}
    \caption{Per-band raw digital number distributions by sensor (log scale; boxes = IQR, whiskers = $1.5\times$IQR, outliers hidden; 400{,}000 sampled pixels).
        Real $\Phi$-sat-2 occupies a scale roughly one order of magnitude below Sentinel-2 and its simulated counterpart, which closely overlap each other.
    }
    \label{fig:band-boxplots}
\end{figure}

Per-domain normalization (Section~\ref{sec:normalization}) closes most, but not all, of this gap (Figure~\ref{fig:band-boxplots-normalized}): after $z$-scoring each sensor independently, the three distributions are centered similarly, but real $\Phi$-sat-2 retains visibly wider spread and longer tails in every band than Sentinel-2 or the simulated domain.
This residual spread difference, not the location/scale difference normalization already removes, is what the domain-adaptation mechanism of Section~\ref{sec:da-mechanism} must close.

\begin{figure}[htp]
    \centering
    \includegraphics[width=\textwidth]{fig_band_boxplots_normalized.png}
    \caption{The same per-band distributions after per-sensor $z$-score normalization of $\sqrt{\text{DN}}$.
        Location and scale are aligned, but real $\Phi$-sat-2 retains a wider spread than Sentinel-2 or the simulated domain in every band: the residual gap normalization alone cannot close.
    }
    \label{fig:band-boxplots-normalized}
\end{figure}

Second, per-band correlation between co-registered real-$\Phi$-sat-2 and Sentinel-2 pixels is only weak-to-moderate even for matched same-scene observations (Pearson $r$ 0.31--0.54, Spearman $\rho$ 0.37--0.61 across the seven shared bands; Figure~\ref{fig:sensor-scatter}), with the weakest correspondence in the red-edge bands (RE1, RE2) and the strongest in the red band.
This is a direct, model-free demonstration that the two sensors disagree substantially about the same physical target observed within the same fortnight.
For contrast, the equivalent cross-band correlation between the simulated domain and its Sentinel-2 source is near-perfect (Spearman $\rho$ 0.98--1.00 on matching bands), as expected since the simulated patches are deterministic transformations of the same Sentinel-2 pixels.
The gap between real and simulated $\Phi$-sat-2 is therefore not explained by any weakness in the correlation analysis itself.

\begin{figure}[htp]
    \centering
    \includegraphics[width=\textwidth]{fig_sensor_scatter.png}
    \caption{Per-band pixel agreement between co-registered real-$\Phi$-sat-2 and Sentinel-2B digital numbers (hexbin density, log scale; orange = ordinary least-squares fit; patches with $\leq$5\% cloud only).
        Point clouds are diffuse around the fitted line in every band, and several bands (RE1, RE2) show a distinct vertical streak near zero $\Phi$-sat-2 DN spanning the full Sentinel-2 DN range, consistent with localized real-sensor dropout or saturation artifacts rather than measurement noise alone.
    }
    \label{fig:sensor-scatter}
\end{figure}

Third, and most directly relevant to the degradation-driver hypothesis discussed in Chapter~\ref{chap:results-discussion}, the shape of the spectral signature (each sensor's per-band median normalized by its own overall band mean, so that absolute scale is divided out) differs between real $\Phi$-sat-2 and Sentinel-2/simulated-$\Phi$-sat-2 (Figure~\ref{fig:spectral-signature}).
Sentinel-2 and the simulated domain are, again, visually indistinguishable from each other.
Real $\Phi$-sat-2, in contrast, shows a pronounced peak at the red band (1.85 relative to its own mean, against 0.85 for Sentinel-2/simulated) and a dip at the near-infrared band (0.70 against 1.21), a difference in relative spectral shape, not only absolute intensity.
This analysis is a quantitative support for attributing the domain gap primarily to spectral response function mismatch (Section~\ref{sec:rq1-results}) rather than to spatial resolution or noise level, since a resolution- or noise-driven gap would not be expected to alter the relative shape of the spectral signature this distinctly.

\begin{figure}[htp]
    \centering
    \includegraphics[width=0.7\textwidth]{fig_spectral_signature.png}
    \caption{Median per-band digital number, normalized by each sensor's own overall band mean (1.0 = that sensor's average).
        Sentinel-2B and simulated $\Phi$-sat-2 overlap almost exactly; real $\Phi$-sat-2 has a markedly different spectral shape, peaking at Red and dipping at RE1 and NIR.
    }
    \label{fig:spectral-signature}
\end{figure}

These measured per-domain statistics (mean, standard deviation, and clip value per band, computed separately for the real, Sentinel-2, and simulated domains) are what parametrize the per-domain normalization protocol of Section~\ref{sec:normalization}.

\subsection{Simulated degraded flood-segmentation dataset}
\label{sec:sen1floods-sim}

No real $\Phi$-sat-2 flood dataset exists, because the mission has not yet captured and labeled a flood event; this is a structural limitation of any new mission, not a data-collection oversight.
To obtain labeled data for the flood-segmentation task motivating this thesis, the Sen1Floods benchmark is passed through the same degradation pipeline used to build the triplets, producing a graded ladder of $\Phi$-sat-2-like flood scenes at increasing degradation severity.
This simulated dataset plays a different role from the triplets as it is the training and evaluation substrate for the flood task itself, not the substrate for measuring the domain gap, since it lacks a real-$\Phi$-sat-2 counterpart against which to validate the simulation directly.

The ladder is produced with an updated implementation of the OrbitalAI simulator, restructured as a configurable sequence of stages inside TerraTorch so that it applies to any Sentinel-2-based dataset in that framework rather than to a single acquisition.
The segmentation mask is carried through the transformation with the imagery (Section~\ref{sec:degradation-pipeline}, Figure~\ref{fig:simulator}).
This reimplementation allows to control each physical effect individually.
This gives the ladder a second purpose beyond supplying a simulated dataset with flood labels.
The triplets establish that a domain gap exists in TerraMind's representation (Section~\ref{sec:gap-measurement}) but not which physical difference between the two instruments produces it, and the two sensors differ in several respects at once, so no measurement on real pairs can separate the causes.
Degrading a fixed set of scenes one factor at a time, at graded severities within each factor makes that separation possible. The scene content is identical across the ladder by construction, so any change in the base TerraMind encoder's response is attributable to the effect being varied rather than to what the imagery depicts.
Chapter~\ref{chap:results-discussion} reports how the teacher responds across this ladder; isolating which individual factor carries the effect is key to answering RQ1.

\section{Physics-Based Sensor Degradation Pipeline}
\label{sec:degradation-pipeline}

The degradation pipeline adapts the simulator developed for the ESA $\Phi$-lab OrbitalAI challenge \cite{longepe2024orbitalai}, which projects labeled Sentinel-2 L1C tiles into the $\Phi$-sat-2 sensor domain by modeling the physical differences between the two instruments as a sequence of operations in radiance space.
Because these effects are defined physically on radiance, not reflectance, the pipeline converts source reflectance to top-of-atmosphere radiance, applies each degradation in turn, and converts back to reflectance to output a simulated L1C product.
Pixel coordinates are written $(i,j)$ and spectral band index $b_n$ throughout.

The implementation used here is an updated version of that simulator, summarised in Figure~\ref{fig:simulator}.
Each stage draws its parameters from its own intermediate input rather than from a single fixed configuration, so an individual effect can be varied or disabled without disturbing the others, which is what permits the factor-isolation use described in Section~\ref{sec:sen1floods-sim}.
Solar geometry and irradiance are retrieved per scene from the Copernicus Data Space OData API using the source product's location and acquisition time, so the radiance conversion of Equation~\eqref{eq:rad2refl} uses the true illumination conditions of each tile rather than a nominal value.

\begin{figure}[htp]
    \centering
    \includegraphics[width=\textwidth]{fig_updated_simulator.png}
    \caption{The updated degradation pipeline.
        A source Sentinel-2 L1 product passes through five physical stages in radiance space before conversion back to reflectance, each driven by its own auxiliary input: the centre-wavelength list and spectral response functions for the spectral adjustment, the 4.75\,m target ground sampling distance for the resampling, a misalignment standard deviation for the band-to-band shift, the per-band effective SNR for the noise term, and the Nyquist MTF for the optical blur.
        Solar zenith angle, Earth--Sun distance, and solar irradiance are retrieved per scene from the Copernicus Data Space OData API.
        Because every stage is parameterised separately, individual effects can be varied or switched off independently, and the segmentation mask is carried through the pipeline so that labels stay aligned with the degraded imagery.
    }
    \label{fig:simulator}
\end{figure}

\paragraph{Radiance--reflectance conversion.} TOA reflectance $\rho$ and TOA radiance $L$ are related by
\begin{equation}
    \rho(i,j,b_n) = \frac{\pi\, d^2}{E_{\text{SUN}}(b_n)\,\cos\theta_S(i,j)}\, L(i,j,b_n),
    \label{eq:rad2refl}
\end{equation}
where $d$ is the Earth--Sun distance in astronomical units, $E_{\text{SUN}}(b_n)$ the exo-atmospheric solar irradiance for band $b_n$, and $\theta_S(i,j)$ the solar zenith angle at pixel $(i,j)$.
The pipeline enters via the inverse of Equation~\eqref{eq:rad2refl} and exits by re-applying it after all radiance-domain degradations; solar irradiance and zenith angle are read from Sentinel-2 tile metadata via the Copernicus Data Space Ecosystem OData API, and no additional solar-irradiance adjustment is needed because Sentinel-2 and $\Phi$-sat-2 spectral bands are close enough in wavelength.
Earth--Sun distance is computed from acquisition day-of-year as $d = 1 - 0.01672\cos(0.9856(day - 4)\cdot \pi/180)$.

\paragraph{Spatial resampling.} $\Phi$-sat-2 has a uniform ground sampling distance of 4.75\,m across all bands; Sentinel-2 provides 10\,m (bands B2, B3, B4, B8) and 20\,m (bands B5, B6, B7).
Each band is resampled to the common 4.75\,m grid using bicubic interpolation, which limits the resampling artifacts that bilinear interpolation would introduce, though it cannot recover spatial detail the coarser source bands never captured.

\paragraph{Band-to-band misalignment.}
Push-broom acquisition co-registers spectral bands imperfectly, since detector rows are exposed at slightly different times and platform micro-vibration perturbs the line of sight during each band's integration window.
For the raw L1A product, the displacement of band $b_n$ relative to the previously acquired band follows
\begin{equation}
    \|\Delta \mathbf{r}(b_n)\| \sim \mathcal{N}(\mu_{b_n}, \sigma^2), \qquad \mu_{b_n} \in [0.44, 1.1]\ \text{pixels},
    \label{eq:misalign-l1a}
\end{equation}
with direction sampled uniformly over $[0, 2\pi)$ and displacement accumulating across preceding bands.
For the geometrically corrected L1C product used throughout this thesis, an on-board co-registration step removes this accumulation, and the residual is instead modeled as a single random displacement of each band relative to a reference band ($\mu_{b_n} = 0$), with standard deviation set higher over water ($\sigma_{\text{sea}} = 6$ pixels) than over land ($\sigma_{\text{land}} = 1$ pixel), reflecting the harder auto-registration problem posed by feature-poor sea surfaces.

\paragraph{Signal-to-noise degradation.} The baseline formulation draws a per-band noise sample from $\mathcal{N}(0,1)$ scaled by the ratio of a reference radiance to a specified SNR,
\begin{equation}
    L_{\text{noise}}(i,j,b_n) = \frac{L_{\text{ref}}(b_n)}{\text{SNR}(b_n)}\,\mathcal{N}(0,1), \label{eq:snr-const} \end{equation} which implicitly treats noise variance as signal-independent, matching the specified SNR only near $L_{\text{ref}}$.
A signal-dependent variant, consistent with a shot-noise regime in which SNR scales with the square root of collected signal, rescales the reference SNR to an effective, radiance-dependent value and samples
\begin{equation}
    L_{\text{noise}}(i,j,b_n) = \frac{\sqrt{L_{\text{in}}(i,j,b_n)\,L_{\text{ref}}(b_n)}}{\text{SNR}(b_n)}\,\mathcal{N}(0,1).
    \label{eq:snr-signal}
\end{equation}
The implementation used for the calibrated baseline degradation level applies Equation~\eqref{eq:snr-signal}.
The extended sweep pushes beyond nominal $\Phi$-sat-2 noise conditions, since domain-invariance has to be stress-tested under more severe degradation than the flight-validated parameters cover.
For that sweep, a simplified variant samples a base SNR uniformly from a specified range and scales noise by the ratio of that SNR to the batch-wise maximum radiance, trading exact photometric fidelity for the ability to sweep degradation severity beyond the calibrated regime.
This trade-off is stated explicitly here because it is a deliberate design choice, not an approximation error: the baseline level is flight-parameter-calibrated, and only the extended sweep sacrifices fidelity for range.

\paragraph{Optical blur via MTF/PSF.}
Optical blur is modeled through the system modulation transfer function (MTF), aggregating static (optics, detector), dynamic (jitter, platform instability), and smearing (motion) contributions.
Assuming shift-invariant blur, imaging reduces to convolution with the system point spread function (PSF).
Given the MTF amplitude at Nyquist frequency $M(b_n)$, the MTF is modeled as Gaussian in normalized spatial frequency $k$,
\begin{equation}
    \text{MTF}(b_n, k) = e^{\ln M(b_n) \cdot k^2},
    \label{eq:mtf}
\end{equation}
with the corresponding PSF obtained as its inverse Fourier transform and convolved with the radiance image.
For $\Phi$-sat-2 the average Nyquist MTF is $M_\Phi(b_n) \approx 0.05$; for Sentinel-2, $M_{\text{S-2}}(b_n) \approx 0.30$.
The reference simulation applies Equation~\eqref{eq:mtf} directly with $M_\Phi(b_n)$, leaving the intrinsic Sentinel-2 blur uncorrected; a combined-filter form exists for the case where source blur should also be compensated, but is not required here since the intrinsic Sentinel-2 blur is small relative to $\Phi$-sat-2's.

This pipeline is shared infrastructure within the research group: it also underlies the pre-training and benchmark datasets used by $\Phi$satNet (Section~\ref{sec:phisatnet-background}).
This thesis uses the pipeline's output in a new way, by stress-testing the base TerraMind encoder across a ladder of degradation severities to isolate which physical effect drives the domain gap, and by using it to produce a simulated flood-segmentation dataset in the absence of any real $\Phi$-sat-2 flood data.

\section{Domain Gap Measurement}
\label{sec:gap-measurement}

\subsection{Normalization protocol}
\label{sec:normalization}

Raw $\Phi$-sat-2 digital numbers are heavy-tailed; a square-root transform is applied first to bring the distribution closer to Gaussian, followed by a smooth (soft-saturating) clip toward a domain-specific clip value rather than a hard clip, and finally per-domain $z$-score standardization using mean and standard deviation statistics computed separately for each sensor domain (real $\Phi$-sat-2, real Sentinel-2, simulated $\Phi$-sat-2), rather than a single shared normalization applied across domains.
An early baseline was corrupted by a normalization defect that caused weight saturation and an artificially depressed performance estimate; correcting it was a precondition for every result reported in Chapter~\ref{chap:results-discussion}.

Because per-domain normalization could itself account for an apparent domain gap that is really just a difference in raw value scale, the domain-gap measurement below was run both with and without per-domain normalization.
The two conditions produced similar overall results, and per-domain normalization produced lower cosine similarity between domains in the encoder's early layers, the opposite of what a normalization artifact would predict, indicating that the measured gap in early layers is a genuine representational effect rather than a calibration artifact.

\subsection{Direct and indirect evidence}

The domain gap is measured directly, in latent space, using the co-registered triplets: paired cosine similarity between the embeddings of the two sensor views of the same ground, proxy $\mathcal{A}$-distance (PAD) computed per encoder layer to locate where the gap concentrates, and a cross-domain linear-probe-transfer test in which a probe fitted on one sensor's embeddings is applied to the other's embeddings of the same patches.
It is measured indirectly through downstream task performance: LULC mIoU of the distilled student relative to a supervised baseline of identical architecture trained directly on the target sensor (Section~\ref{sec:evaluation-protocol}), and flood-segmentation IoU across the simulated degradation ladder of Section~\ref{sec:sen1floods-sim}.

\section{Domain-Adaptation Strategy}
\label{sec:da-mechanism}

\subsection{Fine-tuning on degraded data}
\label{sec:da-negative-result}

The first mechanism attempted was implicit: fine-tune TerraMind's encoder on the LULC task using degraded, simulated training data, on the premise that exposure to degraded examples during training would itself induce a more domain-invariant representation.
Cross-validation against the real-$\Phi$-sat-2/Sentinel-2 triplets showed that it doesn't hold. A model fine-tuned this way doesn't show lower domain gap to real $\Phi$-sat-2 than a model fine-tuned only on clean Sentinel-2.
Standard supervised fine-tuning, even on physics-degraded data produced by a physics-based simulator, doesn't translate into a sensor-invariant representation.

\subsection{Domain-specific batch normalization}
\label{sec:dsbn-mechanism}

The mechanism adopted only requires that the domain of each training example is known.
Domain-Specific Batch Normalization (DSBN) \cite{chang2019dsbn} gives the student a separate batch-normalization affine transform and running statistics per sensor domain, on top of otherwise fully shared convolutional weights.
Concretely, the two-branch distillation scheme (Section~\ref{sec:kd}) forwards a Sentinel-2 batch and a $\Phi$-sat-2 batch through the student separately in each training step; DSBN lets each pass use its own domain's normalization statistics rather than a single statistics set blended across both.
Figure~\ref{fig:dsbn} contrasts the two arrangements.
Under ordinary batch normalization both domains pass through one shared layer, so its running statistics converge to a blend of the two sensors and describe neither exactly.
A DSBN layer instead holds two branches, and each example activates the branch matching its own domain while the inactive branch is left untouched.
Everything outside those layers, which is to say every convolutional weight in the student, stays shared and is therefore trained on both sensors at once, so the mechanism separates the sensor-specific part of the representation (per-channel scale and shift) from the part that ought to be common (the learned filters).

\begin{figure}[htp]
    \centering
    \includegraphics[width=\textwidth]{fig_bn_dsbn.png}
    \caption{Batch normalization compared with domain-specific batch normalization, from Chang et al.~\cite{chang2019dsbn}.
        (a) A single BatchNorm layer processes source ($x_S$) and target ($x_T$) examples alike, so its statistics describe a mixture of the two domains.
        (b) A DSBN layer carries one branch per domain, and each input is routed to the branch corresponding to its own domain; the other branch, shown greyed, is inactive for that example.
        All parameters outside the DSBN layers remain shared across domains.
        In this thesis the source domain is Sentinel-2 and the target domain is $\Phi$-sat-2, and the domain of every input is known both during training and on board.
    }
    \label{fig:dsbn}
\end{figure}
Because the deployed sensor on board $\Phi$-sat-2 is always known at inference, selecting the matching normalization branch at deployment time costs nothing operationally, unlike a sensor-agnostic representation, which would need to be invariant without ever being told which sensor produced a given input.

Since the student must eventually run on the target hardware (Myriad~2), we have to measure the cost of the DSBN mechanism and its implication on the model architecture.
Batch normalization is representable in the deployment format: \texttt{BatchNormInference} belongs to \texttt{opset1}, the earliest OpenVINO operator set, and is therefore within the operator set available to the 2020.3 release that supports this processor.
\footnote{\url{https://docs.openvino.ai/2023.3/openvino_docs_ops_normalization_BatchNormInference_1.html}.
    Operator specifications are versioned by opset rather than by release; the 2020.3 documentation set has been withdrawn by Intel, so the current page is cited for an operator whose opset membership predates it.
}
In practice, however, the operator is not emitted at all, because at inference a batch-normalization layer reduces to a fixed per-channel affine transform.
OpenVINO's model conversion folds the transform into the weights of the preceding convolution by default.
Domain-specific normalization changes only how many such constant sets exist during training, with one per sensor domain.
Because the acquiring instrument is always known on board, only the target-sensor branch is exported, and its statistics fold into the convolution weights as a single-domain model's would.
The deployed graph therefore carries no batch-normalization operation, and DSBN adds neither operator-support risk nor inference cost relative to the supervised baseline: its additional parameters exist only while training.
This argument follows from the documented conversion behaviour of the toolchain rather than from a completed export of the student to the $\Phi$-sat-2 processor, which has not been carried out.

\subsection{Paired-contrastive triplet alignment}
\label{sec:contrastive-mechanism}

The co-registered triplet dataset provides something most domain-adaptation settings lack: for every training example, the identical underlying scene observed by two or three different sensors.
This guarantees a directly supervised alignment objective rather than an unpaired one.
The contrastive loss pulls the embeddings of matched real-$\Phi$-sat-2, real-Sentinel-2, and simulated-$\Phi$-sat-2 patches of the same scene together, while pushing apart embeddings of unrelated scenes, regardless of which sensor produced each embedding.
This directly penalizes the encoder for representing sensor identity, rather than only scene content, and does so using the strongest available supervisory signal, exact scene correspondence, instead of matching marginal distributions between unpaired sensor domains as an MMD-style objective would.

This objective was implemented and evaluated on top of the DSBN configuration above, with the co-registered triplets used as designed: to test a training mechanism, not only to measure the gap.
On the LULC task, it did not improve downstream performance beyond the DSBN-only configuration (Section~\ref{sec:rq2-results}), and all headline results in Chapter~\ref{chap:results-discussion} are reported with its weight set to zero.
This is reported as a negative ablation result rather than omitted: the co-registered pairing that makes a contrastive objective possible is a genuine structural asset of this dataset, and the fact that exploiting it directly added nothing measurable on top of a much simpler per-domain normalization scheme is itself informative about where the sensor gap actually lives (Section~\ref{sec:rq1-results}) and about how much machinery is actually needed to close it.

\section{Cross-Architecture Knowledge Distillation}
\label{sec:kd}

TerraMind's ViT-based encoder cannot run on the $\Phi$-sat-2 on-board processor: the Myriad 2 VPU's supported operator set has no dynamic-axes support, which ViT inference requires regardless of model size.
Therefore, compression must cross architecture families, from ViT teacher to CNN student, rather than compress within the ViT family as some prior on-board GeoFM work does\cite{du2025onorbitdemonstrationgeospatialfoundation}.
The teacher is TerraMind in its base variant, fine-tuned on the Sentinel-2 branch of the triplet dataset for the LULC task and then frozen.
The student is a U-Net\cite{ronneberger2015unet} trained from random initialisation. It was chosen for its simple encoder-decoder structure, which is a natural fit for the segmentation task, and for its low operator set requirements, which are compatible with the Myriad 2 VPU as demonstrated by the $\Phi$satNet baseline (Section~\ref{sec:phisatnet-background}).

Two placements for a domain-adaptation mechanism were considered: baking invariance into the frozen teacher and letting the student inherit it passively by mimicking the teacher's outputs, or giving the student its own adaptation mechanism directly.
This thesis adopts the second, for a reason independent of which specific mechanism is used. The student model, is to be deployed, so whatever mechanism is used should be evaluated on the student's own representation rather than assumed to transfer from the teacher.
Two candidate mechanisms were implemented at the student level: domain-specific batch normalization (Section~\ref{sec:dsbn-mechanism}), adopted as the mechanism responsible for every headline result in Chapter~\ref{chap:results-discussion}, and a paired-contrastive triplet loss (Section~\ref{sec:contrastive-mechanism}), evaluated as an ablation on top of it and found not to improve downstream performance further.
DSBN also incidentally repairs a train/test inconsistency inherent to the two-branch training scheme: the two sensors are forwarded in separate passes, so during training each pass would otherwise be normalized by its own domain's batch statistics while a single shared set of running averages (a blend of both domains) is used at inference, meaning the encoder would never encounter, at test time, the normalization under which it was optimised.
Giving each domain its own running statistics closes this gap as a side effect of the same mechanism.

Figure~\ref{fig:kd-architecture} sets out the resulting training scheme.
Its defining feature is that the teacher is queried only on Sentinel-2, the domain it was pre-trained and fine-tuned on and the only one on which its predictions are trustworthy.
A single teacher forward pass then supervises both student branches, so the student's $\Phi$-sat-2 predictions are distilled against a teacher output computed from the Sentinel-2 view of the same ground.
This transfer of supervision across sensors is possible only because the two views are co-registered, and it keeps the teacher's own domain-shift error out of the student. Querying the teacher directly on $\Phi$-sat-2 would distil that error rather than correct it.

\begin{figure}[htbp]
    \centering
    \resizebox{\textwidth}{!}{%
        \begin{tikzpicture}[
                >=Stealth,
                font=\small,
                database/.style={cylinder, draw, shape border rotate=90, aspect=0.25, align=center,
                        fill=blue!25, minimum height=1.2cm, minimum width=2.8cm, thick},
                frozen/.style={rectangle, draw, rounded corners, align=center, fill=gray!25,
                        minimum width=4.6cm, minimum height=1.4cm, thick},
                model/.style={rectangle, draw, rounded corners, align=center, fill=green!25,
                        minimum width=4.6cm, minimum height=1.6cm, thick},
                feature/.style={rectangle, draw, align=center, fill=yellow!25,
                        minimum width=3.6cm, minimum height=1.0cm, thick},
                labels/.style={rectangle, draw, align=center, fill=blue!12,
                        minimum width=3.2cm, minimum height=1.0cm, thick},
                loss/.style={diamond, draw, align=center, fill=red!20, aspect=2.4, thick,
                        minimum width=2.8cm, inner sep=1pt},
                ablated/.style={diamond, draw=gray, dashed, align=center, fill=black!4,
                        text=gray!45!black, aspect=2.4, thick, minimum width=2.8cm, inner sep=1pt},
                arrow/.style={->, thick},
                backarrow/.style={->, thick, red!75!black, dashed},
            ]

            % --- data -------------------------------------------------------
            \node[database] (data_s2)  at (0,  2.8) {Real Sentinel-2\\(source $x_s$)};
            \node[database] (data_phi) at (0, -1.0) {Real $\Phi$-sat-2\\(target $x_t$)};
            \node[labels]   (mask)     at (0, -4.4) {WorldCover labels\\(shared)};

            % --- models -----------------------------------------------------
            \node[frozen] (teacher) at (6.2,  4.0)
            {TerraMind teacher Base (ViT)\\fine-tuned on Sentinel-2, \textbf{frozen}};
            \node[model]  (student) at (6.2, -0.4)
            {CNN student (U-Net)\\shared convolutional weights\\\textbf{domain-specific BN}};

            % --- outputs ----------------------------------------------------
            \node[feature] (t_soft)   at (11.8,  4.0) {Teacher soft\\predictions $T(x_s)$};
            \node[feature] (s_source) at (11.8,  1.3) {Student logits\\source branch $S(x_s)$};
            \node[feature] (s_target) at (11.8, -1.6) {Student logits\\target branch $S(x_t)$};

            % --- losses -----------------------------------------------------
            \node[loss]    (kd)   at (16.8,  3.0) {KD loss\\KL, $T=4$};
            \node[loss]    (task) at (16.8,  0.2) {Task loss\\focal, class-weighted};
            \node[ablated] (con)  at (16.8, -2.8) {Contrastive terms\\weight $0$ (ablated)};

            % --- forward routing --------------------------------------------
            \draw[arrow] (data_s2.east)  -- ++(1.2,0) |- (teacher.west);
            \draw[arrow] (data_s2.east)  -- ++(1.2,0) |- ([yshift=5mm]student.west);
            \draw[arrow] (data_phi.east) -- ++(1.2,0) |- ([yshift=-5mm]student.west);

            \draw[arrow] (teacher.east) -- (t_soft.west);
            \draw[arrow] (student.east) -- ++(0.9,0) |- (s_source.west);
            \draw[arrow] (student.east) -- ++(0.9,0) |- (s_target.west);

            % one teacher output supervises BOTH student branches
            \draw[arrow] (t_soft.east)   -- ++(1.1,0) |- (kd.north);
            \draw[arrow] (s_source.east) -- ++(0.6,0) |- (kd.west);
            \draw[arrow] (s_target.east) -- ++(1.5,0) |- (kd.south);

            \draw[arrow] (s_source.east) -- ++(0.6,0) |- (task.north);
            \draw[arrow] (s_target.east) -- ++(0.6,0) |- (task.west);
            \draw[arrow] (mask.east) -- (19.9,-4.4) -- (19.9,0.2) -- (task.east);

            \draw[arrow, gray, dashed] (s_target.south) -- ++(0,-1.2) -| (con.west);

            % --- backpropagation --------------------------------------------
            \draw[backarrow] (task.south) -- (16.8,-5.9) -- node[above, black]
            {\footnotesize gradients update the student only; the teacher stays frozen}
            (6.2,-5.9) -- (student.south);

            % --- grouping ---------------------------------------------------
            \begin{scope}[on background layer]
                \node[draw, dashed, thick, blue!45, inner sep=0.35cm, fit=(data_s2) (data_phi) (mask),
                    label=above:\textbf{Co-registered triplet dataset}] {};
                \node[draw, dashed, thick, red!45, inner sep=0.35cm, fit=(kd) (task) (con),
                    label=above:\textbf{Training objective}] {};
            \end{scope}
        \end{tikzpicture}%
    }
    \caption{Cross-architecture distillation with domain-specific normalization.
        The frozen TerraMind teacher is evaluated only on the Sentinel-2 view, where it is trustworthy; its soft predictions supervise both student branches, which is possible because the two sensor views depict the same scene.
        The student is forwarded twice per patch, once per sensor, selecting its own batch-normalization statistics each time while sharing all convolutional weights.
        Both branches also receive the task loss against the single shared WorldCover label map.
        The contrastive terms were implemented and evaluated but carry zero weight in every result reported in Chapter~\ref{chap:results-discussion}.
    }
    \label{fig:kd-architecture}
\end{figure}

The distilled student is therefore domain-adaptive rather than domain-invariant in the strict sense: convolutional weights are shared across sensors, but normalization statistics are not, and the correct branch must be selected at inference according to which sensor produced the input.
This is a weaker property than a single sensor-agnostic representation, but it is the property an on-board deployment actually needs, since the acquiring instrument is always known at inference time.
This is unlike a model expected to generalize to sensors it never saw during training.

\section{Evaluation Protocol}
\label{sec:evaluation-protocol}

Four evaluations follow from the mechanisms above, corresponding to the three research questions of Section~\ref{sec:research-objective}.
First (RQ1), whether the sensor gap is visible in the teacher's representation itself: paired cosine similarity, per-layer PAD, and cross-domain linear-probe transfer between the two sensor views of the same co-registered patches.
Second (RQ2), whether the distilled student's representation inverts these measurements relative to the teacher, and which of the two candidate mechanisms (DSBN, contrastive alignment) is responsible.
Third (RQ3), whether distillation with the adopted mechanism improves LULC task performance over a supervised baseline of identical student architecture trained directly on labelled $\Phi$-sat-2 with no teacher and no second sensor, at a label budget large enough to be representative and, separately, at a much smaller budget, since the direction of this comparison is not assumed to be constant across label budgets.
Fourth, flood segmentation across the simulated degradation ladder, to characterize how gracefully performance degrades with increasing sensor-noise severity.
A few-shot label-efficiency comparison against $\Phi$satNet (Section~\ref{sec:phisatnet-background}) on matched small label budgets, testing directly whether adapting TerraMind's broad, diverse pre-training generalizes to $\Phi$-sat-2 with fewer labels than $\Phi$satNet's mission-specific pre-training requires, remains planned but has not been run (Section~\ref{sec:discussion-comparison}).
The earlier protocol of fully fine-tuning TerraMind directly on real $\Phi$-sat-2 LULC labels, evaluated on the same real domain without any cross-sensor test, produced a baseline mIoU that is retained in Appendix~\ref{app:real-baseline} for reference, but it is not the main baseline of interest because it does not test cross-sensor generalization and therefore does not address the research questions of this thesis.

\section{Tools, Software Environment, and Reproducibility}
\label{sec:reproducibility}

The degradation pipeline and dataset construction are implemented as modular EOTask objects within the \texttt{eo-learn} framework, orchestrated through a versioned simulation configuration object rather than hard-coded parameters, so that a given degradation run (baseline-calibrated or extended-sweep) is fully specified by its configuration rather than by code edits.
Triplet and LULC datasets are stored in HDF5 with an accompanying manifest recording per-patch provenance and train/validation/test split assignment (fixed random seed, positional split), enabling exact reproduction of any reported split.
Model training (TerraMind fine-tuning, DSBN-based distillation, the contrastive-alignment ablation) is implemented in PyTorch.
The codebase is organized as a single versioned Python package with submodules separating data simulation, dataset loading, model training tasks, and drift analysis, under standard git version control, so that each experiment reported in Chapter~\ref{chap:results-discussion} corresponds to a specific, retrievable commit and configuration.
